\section{Complex Dynamics}
\label{sec:mandelbrot}
\index{dynamics!complex}

More generally, we could consider sequences $z_n$ with complex values and thus
recursive equations of the form
\[
  \begin{cases}
    z_0 = \alpha, \\
    z_{n+1} = f(z_n)
  \end{cases}
\]
with $f\colon \CC \to \CC$.
If $f$ were linear, one could explicitly write the solution following the ideas
presented in section~\ref{sec:ricorrenza_lineare}.

We can attempt to consider the simplest nonlinear function that comes to mind,
namely $f(z) = z^2 + c$, and the simplest initial condition $\alpha = 0$. The
problem then becomes studying the behavior of the sequences:
\begin{equation}\label{eq:mandelbrot}
  \begin{cases}
    z_0 = 0, \\
    z_{n+1} = z_n^2 + c
  \end{cases}
\end{equation}
with $c\in \CC$.

The set of real numbers is invariant, thus if $c\in \RR$, the sequence $z_n$
remains real and coincides with the sequence we studied in
exercise~\ref{ex:mandelbrot_reale}.

\begin{figure}
\begin{center}
\includegraphics[width=\textwidth]{mandelbrot.png}
\end{center}
\caption{The Mandelbrot set generated by computer. See the code on page \pageref{code:Mandelbrot}.}
\label{fig:mandelbrot}
\end{figure}

It is very complicated to determine the nature of the sequences defined by
system~\eqref{eq:mandelbrot}.
Only with the use of early computers did Mandelbrot (1924-2010) manage to
graphically represent the \mymargin{Mandelbrot set}%
\index{set!Mandelbrot}%
\index{Mandelbrot!set of}%
set $M$ of points $c\in \CC$ for which the sequence $z_n$ does not diverge: see
figure~\ref{fig:mandelbrot}.

With exercise~\ref{ex:mandelbrot_reale}, we found the intersection $M\cap \RR =
[-2,1/4]$.
In the following exercise, we now aim to demonstrate another simple property
that can be very useful in numerical algorithms used to draw this set.

\begin{exercise}[escape radius]
Demonstrate that if $z_n$ is a solution of \eqref{eq:mandelbrot}, if for some
$N\in \NN$ we have $\abs{z_N}> 2$, then $\abs{z_n}\to +\infty$ as $n\to
+\infty$. In particular, the Mandelbrot set $M$ is contained within the disk
$\ENCLOSE{z\in \CC \colon \abs{z}\le 2}$.
\end{exercise}
%
\begin{proof}[Solution]
Let $c\in \CC$ be fixed and let $z_n$ be the sequence defined
by~\eqref{eq:mandelbrot}.
For every $\eps>0$, consider the set:
\[
    A_\eps = \ENCLOSE{z \in \CC \colon \abs{z}\ge c \text{ and } \abs{z}\ge
  2+\eps}.
\]
We can show that the set $A_\eps$ is invariant because if $z_n\in A_\eps$, we
have
\begin{equation}\label{eq:473244}
\begin{aligned}
\abs{z_{n+1}}
&= \abs{z_n^2+c}
\ge \abs{z_n}^2 - \abs{c}
\ge \abs{z_n}^2 - \abs{z_n}
= (\abs{z_n}-1)\abs{z_n}\\
&\ge (1+\eps)\abs{z_n}
\ge (1+\eps)(2+\eps)
\ge 2+\eps.
\end{aligned}
\end{equation}
Thus, $A_\eps$ is invariant, but not only that, if $z_n \in A_\eps$, we have
found that $\abs{z_{n+1}} \ge (1+\eps) \abs{z_n}$, and thus $\abs{z_{n+k}} \ge
(1+\eps)^k \abs{z_n}$, and therefore $\abs{z_n}\to +\infty$ as $n\to +\infty$.

If $\abs{c}\le 2$ and if $\abs{z_N}> 2$, then $z_N \in A_\eps$ for some
$\eps>0$, and thus we can conclude that $z_n\to \infty\in \bar \CC$.
If instead $\abs{c}>2$, it is easy to observe that $z_0 = 0$, $z_1=c$,
$z_2=c^2+c$, and therefore
\[
  \abs{z_2}
    = \abs{c^2+c}
    \ge \abs{c}^2 - \abs{c}
    = \abs{c}(\abs{c}-1)
    \ge \abs{c} > 2
\]
thus $z_2 \in A_\eps$, and hence, in any case, $z_n\to \infty$.

\end{proof}